% Language: LaTeX
% =====================================================================================================
%  entornos.tex  --  Cajas/entornos propios del documento
%  Requiere: tcolorbox, caption, newfloat (paquetes.tex) e 'ieeeblue' (listings.tex)
% =====================================================================================================
\definecolor{ieeebluebg}{HTML}{EAF1F8}   % fondo claro de la nota IEEE

% --- Float "nota" (estilo figure/table; numeracion "Nota X.Y" + indice) ---
%   \caption dentro numera; \label dentro lo detecta el autocompletado.
%   Indice (opcional):  \listofnota
\DeclareFloatingEnvironment[
  fileext=lonota, listname={Índice de Notas},
  name=Nota, placement=H, within=section]{nota}
\captionsetup[nota]{labelsep=period, font=small, format=hang}

% De querer mantener una nomenclatura de numeros corridos 1..n
% \DeclareFloatingEnvironment[
%   fileext=lonota, listname={Índice de Notas},
%   name=Nota, placement=H]{nota}

% --- Caja de nota IEEE (solo diseno; el caption/label van en el float 'nota') ---
%   Uso simple (sin numerar):   \begin{notaieee}{Titulo} ... \end{notaieee}
%   Uso numerado/referenciable: dentro de \begin{nota} ... \caption{} \label{} \end{nota}
\newtcolorbox{notaieee}[1]{
  enhanced, colback=ieeebluebg, colframe=ieeeblue,
  coltitle=white, colbacktitle=ieeeblue, fonttitle=\bfseries,
  title=#1, boxrule=0.6pt, arc=1.5mm,
  left=2.5mm, right=2.5mm, top=1.2mm, bottom=1.2mm
}